\documentclass{article}
\usepackage[margin=1in]{geometry}
\begin{document}

\section*{Phase 9 Task 3 - Replay, Validation, and Conservative Evaluation}

\textbf{Phase 9 parent doc:} \texttt{Documentation/phases/phase9.tex}\\
\textbf{Task goal:} Run the Phase 5 workflow on the canonical Phase 9 reference hour and produce replay artifacts, holdout validation outputs, conservative paper-trading outputs, and durable summary rows under the canonical \texttt{reports/phase5/} path.

\section*{Why This Task Exists}
Task 2 proved that the chosen reference hour can produce real candidates, evidence, alerts, delivery attempts, and analyst feedback. Task 3 extends that same story into the historical validation layer required by the SRS.
\begin{itemize}
\item The replay bundle must exist for the same exact hour.
\item The validation report must be generated from replay-safe historical timestamps rather than from ``created today'' wall-clock timestamps.
\item Conservative paper-trading outputs must be materialized under \texttt{reports/phase5/}.
\item Matching summary rows must exist in \texttt{replay\_runs}, \texttt{validation\_runs}, and \texttt{backtest\_artifacts}.
\end{itemize}

\section*{Implementation Notes}
\begin{itemize}
\item This task reuses the exact Phase 9 Task 1 window: \texttt{2026-04-20T05:00:00+00:00} to \texttt{2026-04-20T06:00:00+00:00}.
\item It reuses the exact Phase 9 Task 2 fixture candidate chain rather than inventing a second story.
\item Because the native Phase 4 workflow stamps alerts and evidence with the current runtime clock, this task backdates the fixture evidence and alert timeline into the historical hour so Phase 5 can evaluate the window honestly.
\item It also materializes raw-archive coverage, market resolutions, market end dates, and post-alert snapshots for the same fixture markets so the replay and simulator path is non-empty and reviewable.
\end{itemize}

\section*{Canonical Command}
\begin{verbatim}
python run_phase9_phase5_evaluation.py --json
\end{verbatim}

\section*{Outputs}
\begin{itemize}
\item \texttt{reports/phase5/replay\_runs/phase9\_task3/}
\item \texttt{reports/phase5/validation/phase9\_task3\_holdout\_validation.json}
\item \texttt{reports/phase5/validation/phase9\_task3\_holdout\_validation.md}
\item \texttt{reports/phase5/backtests/phase9\_task3\_conservative\_backtest.json}
\item \texttt{reports/phase5/backtests/phase9\_task3\_conservative\_backtest.md}
\item \texttt{reports/phase5/backtests/phase9\_task3\_paper\_trades.json}
\end{itemize}

\section*{Done Condition}
Task 3 is complete when:
\begin{itemize}
\item replay artifacts exist under \texttt{reports/phase5/replay\_runs/},
\item the validation report and conservative backtest outputs exist under \texttt{reports/phase5/},
\item the validation rows for the canonical hour have non-zero coverage and do not depend on future leakage,
\item one new durable row exists in \texttt{validation\_runs},
\item one new durable row exists in \texttt{backtest\_artifacts}, and
\item the resulting report can be cited honestly as a conservative historical evaluation rather than as a planning-only placeholder.
\end{itemize}

\end{document}
